\documentclass[11pt]{article}
\usepackage[a4paper,margin=26mm]{geometry}
\usepackage{amsmath,amssymb,mathtools}
\usepackage{unicode-math}
\setmainfont{Latin Modern Roman}
\setmathfont{Latin Modern Math}
\usepackage{xurl,hyperref}
\hypersetup{hidelinks,pdftitle={Problem statement}}
\newcommand{\sref}[2]{\hyperlink{source:#1:#2}{[S#2]}}
\newcommand{\cataloguescope}[1]{#1}
\begin{document}
\section{Problem 2.10 — (N.}\label{statement}

\subsection{Definitions and mathematical statement}

\paragraph{Shared notation.}$\mathbb N=\{1,2,\ldots\}$, and $\mathbb Z,\mathbb Q,\mathbb R,\mathbb C$ denote the integers, rationals, reals and complex numbers. Any other conventions are specified locally.


A knot or link is a smooth oriented embedding of one or more circles in the oriented $3$-sphere, considered up to orientation-preserving ambient isotopy. Write $O$ for the unknot, which bounds an embedded disk. Let $v_2,v_3$ be the rational primitive finite-type invariants of degrees $2,3$ respectively, normalized by $v_2(T)=v_3(T)=1$ for the right-handed trefoil $T$. Primitive means $v_j(K_1\#K_2)=v_j(K_1)+v_j(K_2)$, where $\#$ is connected sum, and hence $v_j(O)=0$. They obey $v_2(\overline K)=v_2(K)$ and $v_3(\overline K)=-v_3(K)$ for the mirror image $\overline K$; these normalizations specify the invariants in the source.
\paragraph{Statement recorded in the supplied source.}
For crossing counts $n\ge1$, describe the normalized locus
\[\mathcal S=\left\{\left(\frac{v_2(K)}{n^2},\frac{v_3(K)}{n^3}\right):K\text{ admits a diagram with }n\text{ crossings}\right\}\subseteq\mathbb R^2,\]
where $n$ ranges over all positive integers. Include the fixed-$n$ subsets and the limiting region suggested in the source's discussion. Having such a diagram does not assert that $n$ is the minimal crossing number. \sref{statement}{2}
\subsection{Short English statement}
Describe the normalized degree-two and degree-three invariant pairs for knot diagrams with n crossings, including their fixed-crossing and limiting structure.
\subsection{Sources}
\begin{description}
\item[\hypertarget{source:statement:1}{S1}] ULAM. AI and the UnsolvedMath Contributors, \emph{UnsolvedMath: A Curated Collection of Open Mathematics Problems}, catalogue entry. CC BY 4.0. \url{https://huggingface.co/datasets/ulamai/UnsolvedMath}.
\item[\hypertarget{source:statement:2}{S2}] T. Ohtsuki (editor), \emph{Problems on invariants of knots and 3-manifolds}, Geometry \& Topology Monographs 4 (2002), pp. 377--572 (published 2004), item 2.10 and its surrounding definitions and remarks. \url{https://arxiv.org/abs/math/0406190v1}.
\end{description}
\label{end:statement}
\end{document}
